% Draft for ELCVIA (Electronic Letters on Computer Vision and Image Analysis)
% Layout follows the ELCVIA style rules: A4, 11pt Times, single spacing,
% 17 x 25 cm text area, 2 cm top/left margins, figures and tables in the text.
% For the final submission the content can be pasted into the official
% ELCVIA LaTeX template on Overleaf.
\documentclass[a4paper,11pt]{article}
\usepackage[T1]{fontenc}
\usepackage{times}
\usepackage[a4paper,top=2cm,left=2cm,textwidth=17cm,textheight=25cm]{geometry}
\usepackage{graphicx}
\usepackage{booktabs}
\usepackage{amsmath}
\usepackage{xcolor}
\usepackage{url}
\usepackage[hidelinks]{hyperref}
\usepackage{titlesec}
\titleformat*{\section}{\fontsize{14}{16}\bfseries}
\titleformat*{\subsection}{\fontsize{12}{14}\bfseries}
\titleformat*{\subsubsection}{\fontsize{11}{13}\bfseries}
\newcommand{\todo}[1]{\textcolor{red}{[TODO: #1]}}
\setlength{\parskip}{0.3em}

\begin{document}

\begin{center}
{\fontsize{16}{19}\bfseries Does Data Augmentation Close the Robustness Gap Between CNNs and Vision Transformers? A Seen/Unseen Corruption Study on CIFAR-10-C\par}
\vspace{0.8em}
{\large Shynbergen Yegizbayev\par}
\vspace{0.3em}
{\small\em School of Software Engineering, Astana IT University, Astana, Kazakhstan\par}
{\small\em \todo{university e-mail}\par}
\end{center}

\begin{abstract}
\noindent Vision Transformers (ViTs) are often reported to be more robust to common image corruptions than convolutional neural networks (CNNs), and data augmentation is the most common way to improve the robustness of either architecture. It is less clear whether augmentation \emph{closes} the gap between the two families, and whether any reduction holds for corruption types that were not modelled during training. We study this question in a controlled fine-tuning setting. A ResNet-50 and a DeiT-Small, both pretrained only on ImageNet-1k and of similar size (25.6M and 22.1M parameters with the ImageNet head), are fine-tuned on CIFAR-10 with an identical recipe under three conditions: standard augmentation, noise-and-blur corruption augmentation, and AugMix. Robustness is measured on CIFAR-10-C as the relative accuracy drop, separately for corruption families that match the training augmentation (``seen'': noise and blur) and families that do not (``unseen'': weather and digital), with three seeds per setting. With standard training, DeiT-Small loses 12.0\% of its clean accuracy on seen and 4.9\% on unseen corruptions, against 32.3\% and 12.0\% for ResNet-50. Corruption augmentation reduces the absolute gap between the two models by 78\% on seen corruptions (from 20.3 to 4.4 percentage points) but only by 22\% on unseen corruptions (from 7.0 to 5.5 points); measured as a ratio, the unseen gap does not shrink at all. AugMix, which imitates no specific corruption family, reduces the absolute gap by 35\% on seen and 15\% on unseen corruptions, and does not reduce it in ratio terms. The results suggest that augmentation removes the CNN--ViT difference mainly where the test shift was anticipated, while the architectural advantage of the ViT persists on unanticipated shifts.
\end{abstract}

\noindent{\em Key Words}: corruption robustness, vision transformers, convolutional neural networks, data augmentation, distribution shift, CIFAR-10-C.

% =====================================================================
\section{Introduction}

Image classifiers that reach high accuracy on clean test data can fail badly when the input is slightly degraded by noise, blur, weather effects or compression \cite{hendrycks2019benchmarking,geirhos2018generalisation}. Such degradations are common in deployed systems, for example in low-light cameras, moving platforms and compressed video streams. The CIFAR-10-C and ImageNet-C benchmarks \cite{hendrycks2019benchmarking} made this problem measurable by applying 15 synthetic corruptions at five severity levels to standard test sets.

Two lines of work address this problem. The first studies architecture. Since the introduction of Vision Transformers \cite{dosovitskiy2021image,touvron2021training}, several studies have reported that ViTs are more robust to corruptions and other distribution shifts than CNNs \cite{bhojanapalli2021understanding,paul2022vision,naseer2021intriguing}. Later work argued that many of these comparisons were not fair, because the models differed in size and, above all, in their training recipes \cite{bai2021transformers,pinto2022impartial}. Under matched conditions, a modern CNN can be as robust as a Transformer \cite{pinto2022impartial}. The second line studies training data. Data augmentation with noise \cite{rusak2020simple} or with diverse image operations \cite{hendrycks2020augmix} improves corruption robustness considerably. However, the improvement depends strongly on how similar the augmentations are to the test corruptions \cite{mintun2021interaction,yin2019fourier}, and training on one type of distortion does not reliably transfer to others \cite{geirhos2018generalisation}.

These two lines meet in a practical question that, to our knowledge, has received little direct attention: \emph{if a CNN is trained with augmentation, does it catch up with a ViT, and does it catch up also on shifts that the augmentation did not anticipate?} If augmentation closed the gap everywhere, the robustness advantage of ViTs would be mostly a matter of training data. If the gap closes only on the corruption families used in training, the advantage is more likely architectural and remains relevant whenever the deployment shift is unknown.

We answer this question in a small, controlled setting that can be reproduced on free cloud hardware. We fine-tune a ResNet-50 \cite{he2016deep} and a DeiT-Small \cite{touvron2021training} of similar size and identical pretraining data (ImageNet-1k) on CIFAR-10 with the same optimiser, schedule and number of steps. We compare three training conditions and evaluate on CIFAR-10-C with a clear split between corruption families that the augmentation imitates (``seen'') and families it does not (``unseen''). Every setting is repeated with three random seeds.

Our contributions are:
\begin{itemize}
\item A controlled comparison of the CNN--ViT robustness gap under three augmentation conditions, with matched model size, pretraining data and fine-tuning recipe, reported as mean $\pm$ standard deviation over seeds.
\item A seen/unseen analysis showing that noise-and-blur augmentation closes 78\% of the absolute gap on seen corruptions but only 22\% on unseen corruptions, and that the unseen gap does not shrink in relative (ratio) terms.
\item A per-corruption analysis that identifies the corruptions where the gap persists most (pixelation, JPEG compression, contrast and glass blur) and links them to known limits of augmentation transfer.
\end{itemize}

The rest of the paper is organised as follows. Section~\ref{sec:related} reviews related work. Section~\ref{sec:method} describes the models, data, training conditions and metrics. Section~\ref{sec:results} reports the results, Section~\ref{sec:discussion} discusses them together with the limitations of the study, and Section~\ref{sec:conclusion} concludes.

% =====================================================================
\section{Related work}
\label{sec:related}

\subsection{Benchmarks for corruption robustness}
Hendrycks and Dietterich \cite{hendrycks2019benchmarking} introduced ImageNet-C and CIFAR-10-C, which contain 15 corruptions grouped into four families (noise, blur, weather and digital), each at five severities. They also proposed the mean corruption error and a relative error that compares corrupted with clean performance. Earlier, Geirhos et al.\ \cite{geirhos2018generalisation} showed that networks trained on one type of image distortion generalise poorly to other distortions, unlike human observers. Mintun et al.\ \cite{mintun2021interaction} later showed that the benefit of an augmentation is strongly correlated with its perceptual similarity to the test corruptions, and that augmentation gains shrink on corruptions that are dissimilar to those in ImageNet-C.

\subsection{Data augmentation for robustness}
Gaussian noise augmentation, tuned appropriately, is a strong baseline for corruption robustness \cite{rusak2020simple}. Yin et al.\ \cite{yin2019fourier} analysed this effect in the frequency domain: noise augmentation makes models more robust to high-frequency corruptions such as noise and blur, but can make them less robust to low-frequency corruptions such as fog and contrast. AugMix \cite{hendrycks2020augmix} mixes several chains of simple image operations and adds a consistency loss; it deliberately excludes operations that overlap with the corruptions in ImageNet-C and improves robustness across many of them.

\subsection{CNNs versus Vision Transformers}
Vision Transformers \cite{dosovitskiy2021image} split an image into patches and process them with self-attention. DeiT \cite{touvron2021training} showed that ViTs can be trained on ImageNet-1k alone with strong augmentation and regularisation. Several studies then reported that ViTs are more robust than CNNs to corruptions, occlusions and other distribution shifts \cite{bhojanapalli2021understanding,paul2022vision,naseer2021intriguing}. Bai et al.\ \cite{bai2021transformers} argued that earlier comparisons mixed different model sizes and training recipes; with a unified setup, Transformers were not clearly more robust to adversarial attacks, but kept an advantage in out-of-distribution generalisation, which the authors attributed mainly to self-attention. Pinto et al.\ \cite{pinto2022impartial} reached a more balanced conclusion: modern CNNs such as ConvNeXt \cite{liu2022convnet} can be as robust and reliable as Transformers, and both families share similar biases. Wightman et al.\ \cite{wightman2021resnet} showed that ResNet-50 itself improves substantially when trained with a modern, DeiT-like recipe.

\subsection{Position of this work}
Previous comparisons mostly study models trained from scratch on ImageNet, and they treat augmentation either as part of a fixed recipe or as a separate topic. We instead hold the architecture comparison fixed and vary only the fine-tuning augmentation, and we measure the CNN--ViT gap separately on corruption families that the augmentation does and does not imitate. This combines the fairness concerns of \cite{bai2021transformers,pinto2022impartial} with the similarity view of \cite{mintun2021interaction} in a fine-tuning setting, which is how pretrained models are most often used in practice.

% =====================================================================
\section{Method}
\label{sec:method}

\subsection{Research questions}
\begin{itemize}
\item[\textbf{RQ1}] Does data augmentation during fine-tuning reduce the robustness gap between a CNN and a ViT?
\item[\textbf{RQ2}] Does the reduction transfer to corruption types that are not represented in the training augmentation?
\end{itemize}

\subsection{Models}
We compare two models of similar size that were pretrained on the same data (ImageNet-1k, no external data), using the weights distributed with the \texttt{timm} library \cite{wightman2019timm}:
\begin{itemize}
\item \textbf{ResNet-50} \cite{he2016deep} (\texttt{resnet50.a1\_in1k}, 25.6M parameters). These weights were trained with the A1 procedure of \cite{wightman2021resnet}, which uses heavy augmentation and a long schedule similar to the DeiT recipe. This reduces the difference in pretraining recipe between the two models, which was a main criticism of earlier comparisons \cite{bai2021transformers}.
\item \textbf{DeiT-Small} \cite{touvron2021training} (\texttt{deit\_small\_patch16\_224.fb\_in1k}, 22.1M parameters), a ViT with $16\times16$ patches.
\end{itemize}
In both models the ImageNet classification head is replaced by a new linear layer with 10 outputs, and all weights are fine-tuned.

\subsection{Data}
We fine-tune on the 50{,}000 CIFAR-10 training images \cite{krizhevsky2009learning} and evaluate on the 10{,}000 clean test images and on CIFAR-10-C \cite{hendrycks2019benchmarking}. CIFAR-10-C applies 15 corruptions at 5 severities to the same test images, which gives $15\times5\times10{,}000=750{,}000$ test images. All images are $32\times32$ and are resized to $224\times224$ with bilinear interpolation before being passed to the network, so that the input resolution matches pretraining. Corruptions in CIFAR-10-C are applied at $32\times32$, before resizing.

We split the 15 corruptions into two groups according to whether their family is imitated by the corruption augmentation described below (Table~\ref{tab:split}). Note that ``seen'' refers to the \emph{family}: the training augmentation uses our own implementations and parameter ranges, never CIFAR-10-C images.

\begin{table}[t]
\centering
\caption{Split of the 15 CIFAR-10-C corruptions into seen and unseen groups.}
\label{tab:split}
\small
\begin{tabular}{lll}
\toprule
Group & Family & Corruptions \\
\midrule
Seen (7) & Noise & Gaussian noise, shot noise, impulse noise \\
 & Blur & defocus blur, glass blur, motion blur, zoom blur \\
\midrule
Unseen (8) & Weather & snow, frost, fog, brightness \\
 & Digital & contrast, elastic transform, pixelate, JPEG compression \\
\bottomrule
\end{tabular}
\end{table}

\subsection{Training conditions}
All conditions start from the same random crop ($32\times32$ with 4-pixel padding) and horizontal flip, applied at the original resolution.
\begin{enumerate}
\item \textbf{Baseline}: only the crop and flip above.
\item \textbf{Corruption augmentation}: in addition, each training image is corrupted with probability 0.5 by one operation chosen uniformly from five: additive Gaussian noise ($\sigma\sim U[0.02,0.10]$ on the $[0,1]$ scale), Poisson (shot) noise with rate $\lambda\sim U[50,500]$, impulse (salt-and-pepper) noise with amount $\sim U[0.01,0.07]$, Gaussian blur ($5\times5$ kernel, $\sigma\sim U[0.4,1.5]$), and motion blur along a random direction with kernel length in $\{3,5,7\}$. These cover the noise and blur families but none of the weather or digital corruptions.
\item \textbf{AugMix}: in addition, AugMix \cite{hendrycks2020augmix} as implemented in \texttt{torchvision} (severity 3, width 3, random depth). We use the restricted operation set, which excludes contrast, colour, brightness and sharpness operations that overlap with CIFAR-10-C. We apply AugMix as a data transform only, without the Jensen--Shannon consistency loss of the original method, to keep the loss identical across conditions.
\end{enumerate}

\subsection{Fine-tuning recipe}
Both models use the same recipe in all conditions: AdamW \cite{loshchilov2019decoupled} with learning rate $10^{-4}$, weight decay 0.05, batch size 128, 6 epochs (2{,}340 steps), a linear warm-up over the first 5\% of steps followed by cosine decay, label smoothing 0.1, and mixed-precision training. We did not tune hyperparameters separately for each model or condition, so that differences are due to the architecture and the augmentation only. Each setting is run with seeds 0, 1 and 2; the seed controls the head initialisation, the data order and the augmentation randomness. Experiments were run on a single NVIDIA Tesla T4 GPU; one seed takes about 15 minutes of training and 15 minutes of evaluation.

\subsection{Metrics}
Let $A_\text{clean}$ be the clean test accuracy and $A_{c,s}$ the accuracy on corruption $c$ at severity $s$. The relative accuracy drop for corruption $c$ is
\begin{equation}
D_c = \frac{A_\text{clean} - \frac{1}{5}\sum_{s=1}^{5} A_{c,s}}{A_\text{clean}},
\end{equation}
and the drop for a group $G$ (seen or unseen) is the mean of $D_c$ over $c\in G$. A relative measure is used because the two models have different clean accuracy; it is close in spirit to the relative corruption error of \cite{hendrycks2019benchmarking}. We define the \emph{robustness gap} of a training condition as
\begin{equation}
\Delta_G = D_G^\text{ViT} - D_G^\text{CNN},
\end{equation}
in percentage points (pp); a negative value means that the ViT is more robust. Because an absolute difference naturally becomes smaller when both drops become smaller, we also report the ratio $R_G = D_G^\text{CNN}/D_G^\text{ViT}$. All values are reported as mean $\pm$ standard deviation over the three seeds. For the gap we report the standard error of the difference of means and the $p$-value of Welch's $t$-test.

% =====================================================================
\section{Results}
\label{sec:results}

\subsection{Main results}
Table~\ref{tab:main} lists the clean accuracy and the relative drops for all settings, and Figure~\ref{fig:drop} shows the drops on seen and unseen corruptions. Variation across seeds is small: the standard deviation of the clean accuracy is at most 0.13 pp, and that of the group drops at most 1.3 pp.

\begin{table}[t]
\centering
\caption{Clean accuracy and relative accuracy drop on CIFAR-10-C (mean $\pm$ std over 3 seeds). Lower drop is better.}
\label{tab:main}
\small
\begin{tabular}{llcccc}
\toprule
Model & Training & Clean acc. (\%) & Seen drop (\%) & Unseen drop (\%) & All drop (\%) \\
\midrule
ResNet-50 & Baseline & 96.46 $\pm$ 0.09 & 32.25 $\pm$ 1.32 & 11.98 $\pm$ 0.04 & 21.44 $\pm$ 0.63 \\
DeiT-S & Baseline & 98.04 $\pm$ 0.06 & 11.99 $\pm$ 0.52 & 4.94 $\pm$ 0.09 & 8.23 $\pm$ 0.29 \\
ResNet-50 & Corruption aug. & 95.74 $\pm$ 0.13 & 8.38 $\pm$ 0.16 & 8.97 $\pm$ 0.14 & 8.70 $\pm$ 0.06 \\
DeiT-S & Corruption aug. & 98.13 $\pm$ 0.07 & 4.02 $\pm$ 0.12 & 3.52 $\pm$ 0.10 & 3.75 $\pm$ 0.11 \\
ResNet-50 & AugMix & 96.21 $\pm$ 0.04 & 20.35 $\pm$ 0.61 & 9.88 $\pm$ 0.03 & 14.76 $\pm$ 0.29 \\
DeiT-S & AugMix & 98.23 $\pm$ 0.02 & 7.11 $\pm$ 0.51 & 3.90 $\pm$ 0.04 & 5.40 $\pm$ 0.24 \\
\bottomrule
\end{tabular}

\end{table}

\begin{figure}[t]
\centering
\includegraphics[width=\textwidth]{figs/fig_drop_by_condition.pdf}
\caption{Relative accuracy drop on seen (left) and unseen (right) corruptions for each training condition. Error bars show the standard deviation over 3 seeds.}
\label{fig:drop}
\end{figure}

\textbf{Baseline.} DeiT-Small is both more accurate (98.04\% vs.\ 96.46\%) and much more robust than ResNet-50. It loses 12.0\% of its clean accuracy on seen corruptions and 4.9\% on unseen corruptions, compared with 32.3\% and 12.0\% for ResNet-50. The ViT advantage is present for every one of the 15 corruptions (Figure~\ref{fig:percorr}). It is largest for Gaussian, shot and impulse noise, where ResNet-50 loses 43--57\% of its accuracy, and for pixelation, where it loses 33\% against 9\% for DeiT-Small.

\begin{figure}[t]
\centering
\includegraphics[width=\textwidth]{figs/fig_per_corruption_baseline.pdf}
\caption{Relative accuracy drop per corruption for the baseline models (mean over severities and seeds; error bars: std over seeds). The dashed line separates seen and unseen corruptions.}
\label{fig:percorr}
\end{figure}

\subsection{RQ1: does augmentation reduce the gap?}
Corruption augmentation strongly improves ResNet-50 on seen corruptions: the drop falls from 32.3\% to 8.4\%, which is lower than the drop of the baseline DeiT-Small (12.0\%). DeiT-Small also improves, from 12.0\% to 4.0\%. As a result, the absolute gap on seen corruptions shrinks from $-20.3$ to $-4.4$ pp, a 78\% reduction (Table~\ref{tab:gap}). The gap remains statistically significant ($p<10^{-5}$), so augmentation reduces but does not remove it. The cost in clean accuracy is small: $-0.72$ pp for ResNet-50 and $+0.09$ pp for DeiT-Small.

AugMix also helps both models, but less on seen corruptions: the drop of ResNet-50 falls from 32.3\% to 20.3\%, about half of the improvement given by corruption augmentation, and that of DeiT-Small from 12.0\% to 7.1\%. The clean-accuracy cost is small for ResNet-50 ($-0.25$ pp), and DeiT-Small even gains slightly ($+0.19$ pp). The absolute gap on seen corruptions shrinks from $-20.3$ to $-13.2$ pp, a 35\% reduction. In ratio terms, however, AugMix does not reduce the gap: ResNet-50 loses 2.86 times more accuracy than DeiT-Small, compared with 2.69 times at baseline. AugMix thus improves the two architectures by a similar factor, while the targeted corruption augmentation helps the CNN more.

\begin{table}[t]
\centering
\caption{CNN--ViT robustness gap $\Delta_G = D_G^\text{ViT}-D_G^\text{CNN}$ in percentage points ($\pm$ standard error of the difference) and ratio $R_G = D_G^\text{CNN}/D_G^\text{ViT}$. All gaps are significant at $p<0.001$ (Welch's $t$-test).}
\label{tab:gap}
\small
\begin{tabular}{lcccc}
\toprule
& \multicolumn{2}{c}{Seen} & \multicolumn{2}{c}{Unseen} \\
\cmidrule(lr){2-3}\cmidrule(lr){4-5}
Training & Gap (pp) & Ratio & Gap (pp) & Ratio \\
\midrule
Baseline & $-20.26 \pm 0.82$ & 2.69 & $-7.04 \pm 0.06$ & 2.42 \\
Corruption aug. & $-4.36 \pm 0.12$ & 2.09 & $-5.46 \pm 0.10$ & 2.55 \\
AugMix & $-13.24 \pm 0.46$ & 2.86 & $-5.98 \pm 0.03$ & 2.54 \\
\bottomrule
\end{tabular}

\end{table}

\subsection{RQ2: does the reduction transfer to unseen corruptions?}
The transfer to unseen corruptions is weak. With corruption augmentation, the drop of ResNet-50 on unseen corruptions falls from 12.0\% to 9.0\%, and that of DeiT-Small from 4.9\% to 3.5\%. The absolute gap shrinks only from $-7.0$ to $-5.5$ pp (22\%). In relative terms there is no reduction at all: ResNet-50 loses 2.42 times more accuracy than DeiT-Small at baseline and 2.55 times more with corruption augmentation. On seen corruptions, by contrast, the ratio falls from 2.69 to 2.09. AugMix gives a similar, slightly smaller improvement on unseen corruptions (ResNet-50: from 12.0\% to 9.9\%; DeiT-Small: from 4.9\% to 3.9\%), and the absolute gap shrinks from $-7.0$ to $-6.0$ pp (15\%), with the ratio again almost unchanged (2.54). Across all three training conditions the unseen ratio stays between 2.4 and 2.6.

Figure~\ref{fig:severity} shows the same pattern across severities. On seen corruptions the curve of the augmented ResNet-50 is close to, and at high severity above, the baseline DeiT-Small. On unseen corruptions the two ResNet-50 curves stay well below both DeiT-Small curves at every severity, and the difference grows with severity.

\begin{figure}[t]
\centering
\includegraphics[width=\textwidth]{figs/fig_severity.pdf}
\caption{Accuracy as a function of corruption severity, averaged over the seen (left) and unseen (right) corruptions and over 3 seeds.}
\label{fig:severity}
\end{figure}

\subsection{Per-corruption analysis}
Table~\ref{tab:percorr} gives the drop for each corruption. Within the seen group, corruption augmentation almost removes the ResNet-50 weakness on the three noise types (e.g.\ Gaussian noise: from 56.8\% to 9.3\%) but helps much less on glass blur (from 40.6\% to 26.9\%), which is the seen corruption least similar to our Gaussian and motion blur operations, as it combines blur with local pixel shuffling. Within the unseen group, the largest remaining weaknesses of ResNet-50 are pixelation (27.1\% after corruption augmentation, against 6.3\% for DeiT-Small) and JPEG compression (14.6\% vs.\ 8.4\%). Both corruptions remove or distort high-frequency detail in blocks, which the augmentation does not imitate.

\begin{table}[t]
\centering
\caption{Relative accuracy drop (\%) per corruption, mean over severities and 3 seeds. R = ResNet-50, D = DeiT-Small; B = baseline, C = corruption augmentation, A = AugMix.}
\label{tab:percorr}
\small
\begin{tabular}{llcccccc}
\toprule
Group & Corruption & R-B & D-B & R-C & D-C & R-A & D-A \\
\midrule
Seen & Gaussian noise & 56.8 & 19.4 & 9.3 & 4.5 & 40.1 & 13.0 \\
 & Shot noise & 42.8 & 15.0 & 6.7 & 3.2 & 28.7 & 9.5 \\
 & Impulse noise & 48.6 & 12.3 & 4.5 & 2.1 & 21.8 & 3.9 \\
 & Defocus blur & 8.3 & 3.4 & 2.5 & 1.1 & 3.6 & 1.6 \\
 & Glass blur & 40.6 & 22.6 & 26.9 & 13.9 & 35.2 & 16.7 \\
 & Motion blur & 17.3 & 7.0 & 5.5 & 2.0 & 8.1 & 3.1 \\
 & Zoom blur & 11.4 & 4.3 & 3.3 & 1.4 & 4.9 & 2.0 \\
\midrule
Unseen & Snow & 6.5 & 3.4 & 5.5 & 2.6 & 6.3 & 2.9 \\
 & Frost & 8.0 & 3.9 & 5.2 & 2.5 & 7.4 & 3.4 \\
 & Fog & 4.6 & 3.6 & 3.0 & 2.4 & 3.1 & 2.4 \\
 & Brightness & 1.4 & 0.7 & 1.2 & 0.6 & 1.2 & 0.6 \\
 & Contrast & 10.4 & 3.0 & 7.5 & 1.7 & 8.3 & 1.7 \\
 & Elastic transform & 11.8 & 5.6 & 7.7 & 3.7 & 9.0 & 4.3 \\
 & Pixelate & 33.2 & 9.2 & 27.1 & 6.3 & 28.1 & 6.8 \\
 & JPEG compression & 19.9 & 10.2 & 14.6 & 8.4 & 15.6 & 9.0 \\
\bottomrule
\end{tabular}
\end{table}

% =====================================================================
\section{Discussion}
\label{sec:discussion}

\textbf{Augmentation closes the gap where the shift was anticipated.} On noise and blur, a simple augmentation brings ResNet-50 to the robustness level of a ViT trained without it. In this regime, the ``ViT advantage'' can largely be bought with data. This agrees with the view that much of the reported difference between the two families depends on training choices \cite{bai2021transformers,pinto2022impartial}.

\textbf{The architectural advantage persists on unanticipated shifts.} On weather and digital corruptions, the same augmentation improves both models by a similar factor, so the ratio between them stays at about 2.5. A practitioner who does not know the deployment shift in advance therefore still gains robustness by choosing the ViT. This matches the finding of Bai et al.\ \cite{bai2021transformers} that Transformers keep an out-of-distribution advantage under unified training, and the observation of Mintun et al.\ \cite{mintun2021interaction} that augmentation gains depend on similarity to the test corruption. Our contribution is to show both effects in the same controlled experiment and to quantify how much of the gap each explains.

\textbf{Targeted versus generic augmentation.} The two augmentations behave differently. The targeted noise-and-blur augmentation helps the CNN more than the ViT on seen corruptions (the ratio falls from 2.69 to 2.09), because it repairs the CNN's specific weakness to high-frequency noise \cite{yin2019fourier}. AugMix, which does not imitate any CIFAR-10-C family, improves both models by a similar factor on both groups, so it narrows the absolute gap but leaves the ratio unchanged. A generic augmentation therefore does not make a CNN behave more like a ViT; it makes both models more robust.

\textbf{Absolute versus relative gap.} The choice of metric matters. Measured in percentage points, augmentation reduces the unseen gap by 22\%, which could be read as partial transfer. Measured as a ratio, the gap is unchanged. We recommend reporting both, because absolute differences shrink automatically when both models improve.

\textbf{Limitations.} The study has several limitations. (i) It uses one CNN and one ViT of about 22--26M parameters; results may differ for other architectures (e.g.\ ConvNeXt \cite{liu2022convnet}) or larger models. (ii) Only CIFAR-10 is used, with upsampled $32\times32$ images, so resizing may interact with some corruptions, in particular pixelation. (iii) The fine-tuning recipe is short (6 epochs) and was not tuned per model. (iv) AugMix is used without its consistency loss, which is known to contribute to its robustness \cite{hendrycks2020augmix}. (v) Three seeds give stable estimates here, but the standard deviations are themselves uncertain. (vi) ``Seen'' corruptions are matched at the family level, not by exact parameters, so our corruption augmentation is not an oracle for CIFAR-10-C.

% =====================================================================
\section{Conclusion}
\label{sec:conclusion}

We compared the corruption robustness of a ResNet-50 and a DeiT-Small with matched size, pretraining data and fine-tuning recipe, under three augmentation conditions and with a seen/unseen split of CIFAR-10-C. Noise-and-blur augmentation removed 78\% of the absolute CNN--ViT gap on the corruption families it imitates, but only 22\% on other families, and none in relative terms. AugMix narrowed the absolute gap less (35\% and 15\%) and left the ratio unchanged on both groups. Augmentation therefore narrows the gap mainly where the test shift was anticipated, while the ViT keeps an advantage of about 2.5 times smaller accuracy loss on unanticipated shifts. Future work could extend the study to more architectures, to ImageNet-scale data, and to corruption sets that are deliberately dissimilar to the training augmentations.

\section*{Code and data availability}
\todo{GitHub link to the notebook and analysis script.} All experiments use public data (CIFAR-10, CIFAR-10-C) and public pretrained weights.

\section*{Declaration on the use of AI tools}
\todo{Describe honestly how AI assistance (Claude) was used, e.g.\ for code and for language editing, according to the journal's and the course's policy.}

% =====================================================================
\begin{thebibliography}{99}
\small

\bibitem{bai2021transformers}
Y.~Bai, J.~Mei, A.~Yuille, C.~Xie.
\newblock Are Transformers more robust than CNNs?
\newblock In {\em Advances in Neural Information Processing Systems (NeurIPS)}, 2021.

\bibitem{bhojanapalli2021understanding}
S.~Bhojanapalli, A.~Chakrabarti, D.~Glasner, D.~Li, T.~Unterthiner, A.~Veit.
\newblock Understanding robustness of Transformers for image classification.
\newblock In {\em Proceedings of the IEEE/CVF International Conference on Computer Vision (ICCV)}, pages 10231--10241, 2021.

\bibitem{dosovitskiy2021image}
A.~Dosovitskiy, L.~Beyer, A.~Kolesnikov, D.~Weissenborn, X.~Zhai, T.~Unterthiner, M.~Dehghani, M.~Minderer, G.~Heigold, S.~Gelly, J.~Uszkoreit, N.~Houlsby.
\newblock An image is worth 16x16 words: Transformers for image recognition at scale.
\newblock In {\em International Conference on Learning Representations (ICLR)}, 2021.

\bibitem{geirhos2018generalisation}
R.~Geirhos, C.~R.~M.~Temme, J.~Rauber, H.~H.~Sch\"utt, M.~Bethge, F.~A.~Wichmann.
\newblock Generalisation in humans and deep neural networks.
\newblock In {\em Advances in Neural Information Processing Systems (NeurIPS)}, 2018.

\bibitem{he2016deep}
K.~He, X.~Zhang, S.~Ren, J.~Sun.
\newblock Deep residual learning for image recognition.
\newblock In {\em Proceedings of the IEEE Conference on Computer Vision and Pattern Recognition (CVPR)}, pages 770--778, 2016.

\bibitem{hendrycks2019benchmarking}
D.~Hendrycks, T.~Dietterich.
\newblock Benchmarking neural network robustness to common corruptions and perturbations.
\newblock In {\em International Conference on Learning Representations (ICLR)}, 2019.

\bibitem{hendrycks2020augmix}
D.~Hendrycks, N.~Mu, E.~D.~Cubuk, B.~Zoph, J.~Gilmer, B.~Lakshminarayanan.
\newblock AugMix: A simple data processing method to improve robustness and uncertainty.
\newblock In {\em International Conference on Learning Representations (ICLR)}, 2020.

\bibitem{krizhevsky2009learning}
A.~Krizhevsky.
\newblock Learning multiple layers of features from tiny images.
\newblock Technical report, University of Toronto, 2009.

\bibitem{liu2022convnet}
Z.~Liu, H.~Mao, C.-Y.~Wu, C.~Feichtenhofer, T.~Darrell, S.~Xie.
\newblock A ConvNet for the 2020s.
\newblock In {\em Proceedings of the IEEE/CVF Conference on Computer Vision and Pattern Recognition (CVPR)}, pages 11976--11986, 2022.

\bibitem{loshchilov2019decoupled}
I.~Loshchilov, F.~Hutter.
\newblock Decoupled weight decay regularization.
\newblock In {\em International Conference on Learning Representations (ICLR)}, 2019.

\bibitem{mintun2021interaction}
E.~Mintun, A.~Kirillov, S.~Xie.
\newblock On interaction between augmentations and corruptions in natural corruption robustness.
\newblock In {\em Advances in Neural Information Processing Systems (NeurIPS)}, 2021.

\bibitem{naseer2021intriguing}
M.~Naseer, K.~Ranasinghe, S.~Khan, M.~Hayat, F.~S.~Khan, M.-H.~Yang.
\newblock Intriguing properties of vision transformers.
\newblock In {\em Advances in Neural Information Processing Systems (NeurIPS)}, 2021.

\bibitem{paul2022vision}
S.~Paul, P.-Y.~Chen.
\newblock Vision transformers are robust learners.
\newblock In {\em Proceedings of the AAAI Conference on Artificial Intelligence}, 36(2):2071--2081, 2022.

\bibitem{pinto2022impartial}
F.~Pinto, P.~H.~S.~Torr, P.~K.~Dokania.
\newblock An impartial take to the CNN vs Transformer robustness contest.
\newblock In {\em European Conference on Computer Vision (ECCV)}, 2022.

\bibitem{rusak2020simple}
E.~Rusak, L.~Schott, R.~S.~Zimmermann, J.~Bitterwolf, O.~Bringmann, M.~Bethge, W.~Brendel.
\newblock A simple way to make neural networks robust against diverse image corruptions.
\newblock In {\em European Conference on Computer Vision (ECCV)}, 2020.

\bibitem{touvron2021training}
H.~Touvron, M.~Cord, M.~Douze, F.~Massa, A.~Sablayrolles, H.~J\'egou.
\newblock Training data-efficient image transformers \& distillation through attention.
\newblock In {\em International Conference on Machine Learning (ICML)}, pages 10347--10357, 2021.

\bibitem{wightman2019timm}
R.~Wightman.
\newblock PyTorch Image Models.
\newblock \url{https://github.com/huggingface/pytorch-image-models}, 2019.

\bibitem{wightman2021resnet}
R.~Wightman, H.~Touvron, H.~J\'egou.
\newblock ResNet strikes back: An improved training procedure in timm.
\newblock {\em arXiv preprint arXiv:2110.00476}, 2021.

\bibitem{yin2019fourier}
D.~Yin, R.~G.~Lopes, J.~Shlens, E.~D.~Cubuk, J.~Gilmer.
\newblock A Fourier perspective on model robustness in computer vision.
\newblock In {\em Advances in Neural Information Processing Systems (NeurIPS)}, 2019.

\end{thebibliography}

\end{document}
